% !TEX root = ../paper_zh.tex

\section{讨论}

\subsection{从成对证据到场景级几何}

UAV-TAlign-12K 揭示了无人机红外-可见光对齐问题中的一个阶段转变。红外专用方法和现代匹配器
在正式划分的几乎每一对图像上均能生成单应性，从而提供丰富的几何证据。然而，这些证据在空间
支撑、跨模态一致性，以及与同一场景其他帧的一致程度上仍有差异。因此，场景共识是一个独立的
对齐阶段：它将异质逐帧估计转化为可作为统一几何产品使用的场景变换。

消融结果表明，鲁棒共识是贡献最显著的模块。相较仅使用成对证据，共识在提高边缘与梯度一致性的
同时，大幅降低严重离群点比例和变换离散度。跨模态验证随后把稳定几何映射为规范高置信产品集合。
这一组合设计将局部对应关系证据与无人机巡检和多模态感知所需的场景级输出联系起来。

\subsection{覆盖率感知的运行分析}

可靠性--覆盖率曲线为固定阈值评测增加了运行维度。规范运行点提供 9 个高置信产品，并覆盖
74.4\% 的评测图像对；基于分数的扫描则展示更广泛的覆盖区间。应用系统可以据此按照自身质量
要求选择场景产品覆盖率，而无需改变成对匹配器或共识估计器。

两项受控分析支撑了曲线中采用的排序信号。留帧重计算相对完整证据共识仅产生 1.39 像素的漂移
中位数；平移、旋转和尺度扰动也使联合边缘/梯度信号单调下降。这些结果共同把场景诊断量与证据
稳定性和已知几何退化联系起来。

\subsection{基准扩展性}

该基准通过基于 manifest 的接口分离成对对应关系、场景聚合和跨模态验证。未来匹配器可以作为
证据生成器进行比较，也可以在不改变真实采集划分的前提下研究替代共识估计器和运行策略。同一
组织方式还支持扩展到更多场景类别、采集几何、传感器载荷和下游红外感知任务。

确定性证据调度为这一共享接口提供稳定默认策略。多随机种子实验表明，随机选择会改变少量接近
阈值场景的产品状态，但接受帧支撑总体保持稳定。因此，证据调度主要承担可复现性作用，而鲁棒
共识与跨模态验证构成场景产品生成的核心部分。
